\chapter{The status of the headline results}\label{ch:statusall}

Table~\ref{tab:statusall} gathers, in one place and in the order of the seven parts, the results a referee is likely to check first.
Each row gives the value as it stands in its chapter, the status label it carries there, and the chapter that establishes it; the row on Euclid cites the published forecasts. The
labels are those of the key in the front matter: \derived{} follows from stated premises with no parameter adjusted to data;
\calc{} is computed from a derived formula and published constants or measurements; \calibrated{} is set once from reference data and
then held fixed; \measured{} is read from data by the instruments of this book; \observed{} is a published measurement; \prediction{}
names a measurement not yet made; \conjecture{} is a claim that is not yet a calculation; \openprob{} is a known gap. A row with two
labels carries both, in the order of its parts.
The predictions themselves, with their test dates and falsifiers, are listed in Chapter~\ref{ch:predictions}.

{\small\setlength{\tabcolsep}{3pt}\begin{longtable}{@{}>{\raggedright\arraybackslash}p{0.28\textwidth}>{\raggedright\arraybackslash}p{0.29\textwidth}>{\raggedright\arraybackslash}p{0.13\textwidth}>{\raggedright\arraybackslash}p{0.255\textwidth}@{}}
\caption{Status of the headline results across the seven parts.}\label{tab:statusall}\\
\toprule result & value & status & where \\\midrule\endfirsthead
\toprule result & value & status & where \\\midrule\endhead
\bottomrule\endfoot
```latex
\multicolumn{4}{@{}l}{\emph{Part~\ref{part:1}: IAM's Law and the virial theorem}}\\\cmidrule[0.2pt]{1-4}
Cost per bit from the cell to the cosmic horizon & falls by $1.2\times10^{32}$ & \calc & Ch.~\ref{ch:surfaces} \\\cmidrule[0.2pt]{1-4}
Entropy-area coefficient & $\eta=c^3/(4\hbar G)=1/(4\ell_P^2)$, $\tfrac14=2\pi/8\pi$; one nat per $4\ell_P^2$ & \derived{} (given $G$) & Ch.~\ref{ch:iams_law} \\\cmidrule[0.2pt]{1-4}
The cell's CpG sites against the holographic capacity of a nucleus-sized area & $2.82\times10^{7}$ sites, 52 orders of magnitude below $1.6\times10^{59}$ bits & \calc & Ch.~\ref{ch:surfaces} \\\cmidrule[0.2pt]{1-4}
Equation of state of the record term & $w_{\rm info}=-1-1/(3a)$; $-1.33$ today & \derived & Ch.~\ref{ch:iams_law} \\\cmidrule[0.2pt]{1-4}
Drive in thermal units for a black hole, $Mc^2/k_BT_{\rm BH}=2S_{\rm BH}/k_B$ & $2.1\times10^{77}$ at one solar mass & \calc & Ch.~\ref{ch:iams_law} \\\cmidrule[0.2pt]{1-4}
Virial share in hydrogen & $\langle K\rangle/|\langle V\rangle|=13.606/27.211=\tfrac12$ & \derived & Ch.~\ref{ch:virial_law} \\\cmidrule[0.2pt]{1-4}
Smarr share of Schwarzschild horizons, 1 to $6.5\times10^9\,M_\odot$ & $T_HS/Mc^2=0.5000000000$ & \calc & Ch.~\ref{ch:virial_law} \\\cmidrule[0.2pt]{1-4}
Virial ratio of simulated halos & $2K/|W|=1.1$--$1.3$ (1.02--1.17 with surface pressure) & \observed & Ch.~\ref{ch:virial_law} \\\cmidrule[0.2pt]{1-4}
The cosmic horizon at one half & $\beta_m/\Omega_m=1/2$ & \prediction & Ch.~\ref{ch:virial_law} \\\cmidrule[0.2pt]{1-4}
\multicolumn{4}{@{}l}{\emph{Part~\ref{part:2}: cosmology}}\\\cmidrule[0.2pt]{1-4}
Deformed-entropy cosmologies (Barrow, Tsallis) against an added source & their exponent $\Delta$ or $\delta$ is fitted; IAM keeps the area law and fixes $\beta_m$ from $\Omega_m$ & \interp & Ch.~\ref{ch:theory} \\\cmidrule[0.2pt]{1-4}
The cosmological coupling, fixed in every chain & $\beta_m=\Omega_m/2=0.15765$ & \prediction & Ch.~\ref{ch:dual} \\\cmidrule[0.2pt]{1-4}
Coupling left free against the local distance ladder and $f\sigma_8$ & $0.155\pm0.029$, the local rate restated; $f\sigma_8$ alone allow $-0.28\le\beta\le0.47$ & \calc & Ch.~\ref{ch:dual} \\\cmidrule[0.2pt]{1-4}
Collapse exponent required by $E(a)=e^{1-1/a}$ & $n=7/2$ & \calc & Ch.~\ref{ch:theory} \\\cmidrule[0.2pt]{1-4}
Dark-sector equation of state, CPL tangent at $a=1$ & $w_0=-1.062$, $w_a=-0.012$ & \derived & Ch.~\ref{ch:theory} \\\cmidrule[0.2pt]{1-4}
Lensing response & $\Sigma=1$ at every redshift & \calc & Ch.~\ref{ch:theory} \\\cmidrule[0.2pt]{1-4}
Matter-sector coupling today & $\mu_0=-0.136$; $1-\mu=13.62\,\%$ (not the growth deficit) & \calc & Ch.~\ref{ch:surveys} \\\cmidrule[0.2pt]{1-4}
Level~2 fit to Planck 2018, coupling fixed & $\Delta\chi^2=+0.54$ (IAM higher), no added parameter & \measured & Ch.~\ref{ch:level2} \\\cmidrule[0.2pt]{1-4}
Level~1 fits, coupling fixed & $\Delta\chi^2=+0.96$ (Planck), $+0.56$ (+RSD), $+1.73$ (+BAO), $+1.58$ (+Pantheon+) & \measured & Ch.~\ref{ch:latetime} \\\cmidrule[0.2pt]{1-4}
$\sigma_8$, Level~2 (Run~C $\to$ Run~A) & $0.8087\to0.7998$ ($-1.1\,\%$, $-1.51\sigma$) & \measured & Ch.~\ref{ch:level2} \\\cmidrule[0.2pt]{1-4}
$S_8$, Level~2 & $0.830\to0.822$ ($-0.78\sigma$) & \measured & Ch.~\ref{ch:level2} \\\cmidrule[0.2pt]{1-4}
Every other Level~2 parameter & moves by under $0.1\sigma$ & \observed & Ch.~\ref{ch:level2} \\\cmidrule[0.2pt]{1-4}
$\sigma_8$, Level~1, each data combination & $-1.6\,\%$ & \measured & Ch.~\ref{ch:latetime} \\\cmidrule[0.2pt]{1-4}
Amplitude $\mu_0$ left free & medians $+0.030$ to $+0.064$; posterior reaches the $+0.2$ prior edge; prediction inside every 90\,\% interval & \calc & Ch.~\ref{ch:latetime} \\\cmidrule[0.2pt]{1-4}
Photon-sector Hubble constant (Level~2 chain) & $67.16\pm0.47$ km\,s$^{-1}$\,Mpc$^{-1}$; $-0.37\sigma$ from Planck & \measured & Ch.~\ref{ch:dual} \\\cmidrule[0.2pt]{1-4}
Matter-sector Hubble constant, $H_0\sqrt{1+\beta_m}$ & $72.26\pm0.50$; $-0.75\sigma$ from SH0ES & \derived{} (relation), \calc{} (value) & Ch.~\ref{ch:dual} \\\cmidrule[0.2pt]{1-4}
The term placed in the background Friedmann equation & $H_0=61.45\pm0.42$, $10.9\sigma$ from Planck: background form excluded, the large distortion the test was built to show; in IAM the background does not change & \measured & Ch.~\ref{ch:level2} \\\cmidrule[0.2pt]{1-4}
Supernova Hubble-diagram shape (Pantheon+SH0ES, 1590 SNe) & $\beta_{\rm distance}=-0.035$ (68\,\%: $-0.065$ to $0.000$): $\Lambda$CDM geometry & \observed & Ch.~\ref{ch:dsvalidation} \\\cmidrule[0.2pt]{1-4}
Supernova distances with the $\beta$ term applied & $\beta$ on distances excluded, $\Delta\chi^2=+23.6$: distances follow the photon sector, as the dual-sector picture predicts & \calc & Ch.~\ref{ch:dsvalidation} \\\cmidrule[0.2pt]{1-4}
Effective growth index today & $\gamma=0.585$ ($\Lambda$CDM 0.554; measured $0.633^{+0.025}_{-0.024}$) & \calc, \observed & Ch.~\ref{ch:s8trend} \\\cmidrule[0.2pt]{1-4}
$f\sigma_8$ deficit against $\Lambda$CDM, same early amplitude & 4.25\,\% ($z=0$), 2.17\,\% (0.3), 1.35\,\% (0.5), 0.41\,\% (1) & \prediction & Ch.~\ref{ch:surveys} \\\cmidrule[0.2pt]{1-4}
$E_G$ statistic & $+1.8\,\%$ at $z=0.3$, $+3.6\,\%$ today & \calc & Ch.~\ref{ch:surveys} \\\cmidrule[0.2pt]{1-4}
CMB lensing power & $C_L^{\phi\phi}$ lower by 0.08\,\% & \calc & Ch.~\ref{ch:surveys} \\\cmidrule[0.2pt]{1-4}
ISW--galaxy cross-correlation amplitude & $A_{\rm ISW}^{\rm IAM}/A_{\rm ISW}^{\Lambda\rm CDM}\approx1.03$ & \calc & Ch.~\ref{ch:surveys} \\\cmidrule[0.2pt]{1-4}
Euclid test of the IAM $\mu(z)$ & published errors on $1+\mu_0$: 23.3\,\% (conservative), 4\,\% (weak lensing and photometric clustering to $k\approx4\,{\rm Mpc}^{-1}$), about 1\,\% (optimistic)~\cite{Albuquerque2025,Frusciante2025}; for the IAM $\mu(z)$ about $0.3\sigma$ to about $7\sigma$. A forecast with the IAM $\mu(z)$ itself is open & \observed, \calc, \openprob & Ch.~\ref{ch:latetime} \\\cmidrule[0.2pt]{1-4}
Siren population needed to separate the two Hubble constants at $3\sigma$ & $\sigma(H_0)\le1.70$ km\,s$^{-1}$\,Mpc$^{-1}$ & \calc & Ch.~\ref{ch:surveys} \\\cmidrule[0.2pt]{1-4}
Level~2 $\sigma_8$ against lensing surveys & joint KiDS-Legacy + DES Y3 + DESI $0.1\sigma$; KiDS-Legacy $S_8=0.815$, $0.3\sigma$; DES Y3 3$\times$2pt 0.776, $2.3\sigma$ & \observed & Ch.~\ref{ch:sectortension} \\\cmidrule[0.2pt]{1-4}
CPL image of $w_{\rm info}$ against DESI DR2 & $(-\tfrac43,-\tfrac13)$, 7.6--10.2$\sigma$ from DESI in $w_0$; DESI distances are on the photon ruler, so this is not a test & \calc & Ch.~\ref{ch:wzfuture} \\\cmidrule[0.2pt]{1-4}
CONFLICT
Asymptotic expansion rates & 55.78 (light), 71.12 (matter) vs. 55.57 (light), 70.86 (matter) at $H_0=67.16$; no Big Rip & \calc, \derived & Ch.~\ref{ch:wzfuture} \\\cmidrule[0.2pt]{1-4}
Ratio of matter to photon rate & $H_m/H=1.076$ today, 1.275 as $a\to\infty$ & \calc & Ch.~\ref{ch:darkenergy} \\\cmidrule[0.2pt]{1-4}
Writing rate of the record & peak 3.38\,\%\,Gyr$^{-1}$ at $z=1.26$; 2.54\,\%\,Gyr$^{-1}$ now & \calc & Ch.~\ref{ch:darkenergy} \\\cmidrule[0.2pt]{1-4}
Baryon fraction relation $\Omega_b/\Omega_m\approx(3/16)\sqrt{\Omega_\Lambda}$ & holds to 0.79\,\%; derivation of 3/16 open & \observed, \openprob & Ch.~\ref{ch:lambda} \\\cmidrule[0.2pt]{1-4}
Baryon-to-photon ratio from the CMB alone & $\eta=(6.113\pm0.037)\times10^{-10}$ & \observed & Ch.~\ref{ch:baryon} \\\cmidrule[0.2pt]{1-4}
Lensing-to-dynamical mass ratio, Level~1 form only & $1/\mu(z)$; whether it applies in virialised clusters is open & \calc, \openprob & Ch.~\ref{ch:lensdyn} \\\cmidrule[0.2pt]{1-4}
\multicolumn{4}{@{}l}{\emph{Part~\ref{part:bh}: black holes and horizons}}\\\cmidrule[0.2pt]{1-4}
Kerr horizons: bits at the horizon temperature & $T_{\rm BH}S_{\rm BH}/Mc^2=\sqrt{1-\chi^2}/2$ & \calc & Ch.~\ref{ch:blackholes} \\\cmidrule[0.2pt]{1-4}
Stefan--Boltzmann luminosity of a horizon against the Hawking luminosity & $P_{\rm SB}/P_{\rm H}=1$ & \calc & Ch.~\ref{ch:blackholes} \\\cmidrule[0.2pt]{1-4}
Black-body evaporation time of one solar mass & $2.1\times10^{67}$\,yr & \calc & Ch.~\ref{ch:blackholes} \\\cmidrule[0.2pt]{1-4}
\multicolumn{4}{@{}l}{\emph{Part~\ref{part:qp}: particles and quantum records}}\\\cmidrule[0.2pt]{1-4}
Bottom-up collapse exponent from halo mass functions & passes $7/2$ at $z\approx3$--4 & \calc & Ch.~\ref{ch:quantumrecords} \\\cmidrule[0.2pt]{1-4}
Koide ratio of the charged leptons & $Q=0.66
\end{longtable}}

Four points of reading follow from the table. First, the fits to data in Part~\ref{part:2} are made with the coupling fixed by the
law, so the chain values are measurements of the fixed model, not fits of a free parameter; the free-amplitude runs are reported as
medians and bounds because their posteriors reach the prior edge. Second, every floor of the processors in Part~\ref{part:3} is the law at the temperature the record sees, so a device reading needs no
fitted constant; what remains open there is how close each loss channel comes to its floor. Third, the cell gauge of
Part~\ref{part:4} is commissioned on its healthy reference; where breach and the cancer region lie on it is still to be measured.
Fourth, ``no free parameter'' has a precise meaning here. No parameter is fitted to the data the law is tested against: $\beta_m$ comes from
$\Omega_m$; the cell's references are read from healthy cells and then held fixed; the chains with $\mu_0$ free add nothing to the model. That is
narrower than ``no free parameter'' without qualification. Counted in full, the results use calibrated quantities (the cell references held fixed after one
measurement), measured ones (the holding energy $E_{\rm hold}$ and the position $P$), and one factor found numerically, $(2\pi)^{3/10}$ in the electron
mass; the $3/16$ of the baryon relation and the $2/\pi$ of the vacuum estimate are taken as given here; their derivation is an open problem. Each is listed with its label in the rows above. \interp

Every row carries its label, so a reader always knows how strongly each result is claimed and where the next test would bite. Pick a row in your field and push on it.
